% ==========================================
% BAB I PENDAHULUAN
% ==========================================
\cleardoublepage
\chapter{PENDAHULUAN}
\label{chap:pendahuluan}

\section{Latar Belakang}

Kelancaran arus kontainer berpengaruh langsung terhadap biaya dan kualitas
layanan logistik pelabuhan. Di Indonesia, biaya logistik terhadap Produk
Domestik Bruto (PDB) tercatat pernah mencapai 24\% pada tahun 2020 dan menurun
menjadi sekitar 14,29\% pada tahun 2023; angka tersebut masih berada di atas
pembanding kawasan seperti Singapura yang berada pada kisaran 8\%
\autocite{kemenkeu2023nle,pelindo2023logistik}. Salah satu proses yang
mendukung kelancaran arus tersebut adalah inspeksi kontainer, karena hasil
inspeksi menentukan kesiapan fisik kontainer, kelengkapan dokumen, dan tindak
lanjut operasional.

Kebutuhan inspeksi semakin penting karena kondisi fisik kontainer tidak selalu
seragam. Survei Kementerian Perhubungan Indonesia pada tahun 2018 menunjukkan
bahwa hanya 20\% kontainer di Indonesia yang layak pakai, sedangkan 80\%
lainnya berada dalam kondisi buruk atau tidak memenuhi standar kualitas
\autocite{lukmandono2024}. Namun, peningkatan proses inspeksi tidak cukup hanya
dilihat sebagai persoalan alat pemeriksaan. Hasil identifikasi kontainer, bukti
visual, status pemeriksaan, dan penanggung jawab inspeksi perlu terhubung dalam
riwayat kontainer yang konsisten agar dapat digunakan kembali untuk pelacakan
dan audit.

Hasil survei lapangan yang dirangkum pada Lampiran~\ref{chap:survei}
menunjukkan bahwa proses eksisting masih melibatkan pemeriksaan visual,
dokumen kertas, pencatatan lapangan, dan koordinasi lintas pihak. Sistem
operasional terminal dapat mengelola pergerakan kontainer, tetapi data hasil
inspeksi belum selalu tersambung secara otomatis dengan identitas kontainer dan
bukti visual. Akibatnya, informasi inspeksi berisiko tersebar pada dokumen,
catatan petugas, atau sistem yang berbeda sehingga riwayat kontainer tidak
terbentuk secara utuh.

Permasalahan tersebut menunjukkan perlunya rancangan arsitektur sistem yang
menyatukan akuisisi data di lapangan, pemrosesan awal, penyimpanan hasil
inspeksi, dan penyajian riwayat kontainer. Peraturan Menteri Perhubungan Nomor
PM~8 Tahun 2022 mengatur pelayanan kapal melalui Inaportnet sebagai layanan
elektronik di pelabuhan \autocite{kemenhub2022}. Oleh karena itu, tugas akhir
ini berfokus pada perancangan arsitektur sistem otomasi inspeksi kontainer di
pelabuhan, dengan kontainer sebagai pusat informasi dan inspeksi sebagai
aktivitas yang membentuk riwayat kontainer.

\section{Rumusan Masalah}

Berdasarkan latar belakang yang telah diuraikan, permasalahan utama yang
menjadi fokus tugas akhir ini dirumuskan sebagai berikut:
\begin{enumerate}
  \item Bagaimana merancang arsitektur sistem yang mengintegrasikan
  komponen \textit{edge}, layanan aplikasi, penyimpanan data, dan antarmuka
  pengguna dalam satu alur kerja yang konsisten?
  \item Bagaimana merancang struktur data pemeriksaan kontainer agar pencatatan
  hasil inspeksi, riwayat kontainer, dan penyediaan bukti inspeksi dapat
  dilakukan secara terstruktur?
  \item Bagaimana merancang alur informasi yang memungkinkan hasil
  identifikasi dan hasil pemindaian dari beberapa komponen dapat dikaitkan pada
  entitas kontainer yang sama?
  \item Bagaimana merancang pembagian fungsi antarkomponen agar pemrosesan
  lokal, persistensi data, dan penyajian informasi dapat berjalan selaras dengan
  kebutuhan operasional inspeksi kontainer?
\end{enumerate}

Untuk menjawab rumusan masalah tersebut, tugas akhir ini terlebih dahulu
melakukan identifikasi kebutuhan sistem sebagai dasar perancangan, kemudian
memeriksa kesesuaian rancangan terhadap sistem yang direalisasikan sebagai
bukti bahwa rancangan dapat diwujudkan dalam ruang lingkup tugas akhir.

\section{Tujuan}

Tugas akhir ini bertujuan menghasilkan rancangan arsitektur sistem otomasi
inspeksi kontainer yang dapat menjadi dasar realisasi sistem dan
menjadi acuan pengembangan lebih lanjut di lingkungan pelabuhan.

Secara khusus, tugas akhir ini memiliki tujuan sebagai berikut:
\begin{enumerate}
  \item Merancang arsitektur sistem yang mendefinisikan komponen utama,
  hubungan antarkomponen, dan alur informasi inspeksi kontainer secara
  menyeluruh.
  \item Merancang arsitektur data yang mendukung pencatatan hasil inspeksi,
  penyimpanan bukti visual, dan riwayat pemeriksaan kontainer.
  \item Menyusun model alur informasi yang memungkinkan integrasi hasil
  identifikasi, hasil pemindaian, dan penyajian informasi kepada pengguna dalam
  satu sistem yang konsisten.
  \item Merancang pembagian fungsi antarkomponen agar pemrosesan lokal,
  persistensi data, dan penyajian informasi dapat berjalan sesuai kebutuhan
  operasional inspeksi kontainer.
\end{enumerate}

\section{Batasan Masalah}

Agar pembahasan tetap fokus pada kontribusi arsitektur sistem, ruang lingkup
tugas akhir ini dibatasi sebagai berikut:
\begin{enumerate}
  \item Tugas akhir ini berfokus pada perancangan arsitektur sistem dan
  arsitektur data, bukan pada pengembangan detail algoritma deteksi atau
  perangkat keras sensor.
  \item Rancangan arsitektur dibatasi pada definisi komponen utama, relasi
  antarkomponen, struktur data, dan alur informasi yang mendukung proses
  otomasi inspeksi kontainer.
  \item Tugas akhir ini tidak membahas rancangan infrastruktur fisik pelabuhan,
  operasi bisnis pelabuhan secara menyeluruh, atau perancangan ulang seluruh
  ekosistem sistem informasi pelabuhan.
  \item Evaluasi tugas akhir ini difokuskan pada kesesuaian rancangan arsitektur
  terhadap sistem yang direalisasikan dan pada bukti integrasi antarkomponen,
  bukan pada uji performa kuantitatif menyeluruh atau uji
  lapangan skala penuh.
  \item Aspek keamanan siber, tata kelola layanan, dan pengelolaan operasional
  dibahas pada batas keputusan arsitektur, yaitu kontrol akses, pencatatan
  aktivitas, pemisahan tanggung jawab komponen, dan kebutuhan integrasi.
\end{enumerate}

\section{Definisi Operasional Istilah}

Untuk menghindari perbedaan pemaknaan istilah, beberapa istilah utama dalam
tugas akhir ini digunakan dengan makna operasional berikut:
\begin{itemize}[leftmargin=1.2cm,itemsep=0.15em,topsep=0.2em]
  \item \textbf{Kontainer}: unit fisik standar yang menjadi objek utama inspeksi, identifikasi, pencatatan hasil, dan pembentukan riwayat pemeriksaan.
  \item \textbf{Inspeksi kontainer}: pemeriksaan terhadap kontainer yang mencakup identifikasi nomor, pemeriksaan visual, pencatatan hasil, dan pengaitan bukti ke entitas kontainer.
  \item \textbf{Manifes}: data deklaratif mengenai muatan dan perjalanan kontainer yang digunakan sebagai konteks pembanding terhadap hasil identifikasi atau pemindaian.
  \item \textbf{Riwayat kontainer}: kumpulan informasi yang dapat ditelusuri mengenai identitas, waktu dan status pemeriksaan, hasil inspeksi, penanggung jawab, serta referensi bukti visual.
  \item \textbf{Artefak visual}: berkas citra atau video hasil akuisisi dan pemindaian yang menjadi bukti pendukung proses inspeksi.
\end{itemize}

\section{Pembagian Tugas Kelompok}

Tugas akhir ini dikerjakan dalam konteks proyek kelompok Cargovision dengan
ruang lingkup yang saling melengkapi. Untuk memperjelas batas kontribusi, Tabel
\ref{tbl:pembagian-tugas-kelompok} merangkum fokus pengerjaan setiap anggota.
Laporan ini berfokus pada bagian arsitektur sistem, kontrak data berbasis
kontainer, dan batas tanggung jawab antarkomponen.

\begin{table}[ht]
  \centering
  \small
  \caption[Pembagian tugas kelompok Cargovision]{Pembagian Tugas Kelompok Cargovision}
  \label{tbl:pembagian-tugas-kelompok}
  \setlength{\tabcolsep}{4pt}
  \renewcommand{\arraystretch}{1.12}
  \begin{tabularx}{\textwidth}{>{\raggedright\arraybackslash}p{0.28\textwidth} >{\raggedright\arraybackslash}X >{\raggedright\arraybackslash}p{0.24\textwidth}}
    \toprule
    \textbf{Anggota} & \textbf{Fokus Kontribusi} & \textbf{Keterkaitan dengan Laporan Ini} \\
    \midrule
    Aththariq Lisan Quran Daulah Sentono & Perancangan arsitektur, kontrak data berpusat pada kontainer, serta batas tanggung jawab antara \textit{edge}, API, web, dan penyimpanan. & Menjadi fokus utama tugas akhir ini. \\
    Alessandro Jusack Hasian & Realisasi prototipe, alur OCR dan YOLO11, layanan \textit{edge}, backend, dasbor web, dan alur peninjauan operator. & Menjadi artefak realisasi yang digunakan sebagai bukti evaluasi arsitektur. \\
    Eleanor Cordelia & Tata kelola, COBIT/ITIL/SLA, jejak audit, \textit{API gateway}, \textit{anti-corruption layer}, serta model peran dan akuntabilitas. & Menjadi konteks pendukung untuk aspek kontrol akses, audit, dan kesiapan tata kelola. \\
    Muhammad Faiz Atharrahman & Model penempatan, rancangan hibrida \textit{on-premise} tiga lapis, CI/CD, VPN, pemantauan, dan pengujian beban. & Menjadi konteks pendukung untuk aspek penempatan, kapasitas, dan evaluasi operasional lanjutan. \\
    \bottomrule
  \end{tabularx}
\end{table}

\FloatBarrier

\section{Metodologi}

Tugas akhir ini menggunakan pendekatan \textit{Design Science Research} (DSR)
karena hasil utama yang disusun adalah artefak rancangan arsitektur sistem
\autocite{hevner2004design,peffers2007dsrm}. Pendekatan ini digunakan untuk
menghubungkan masalah operasional di lapangan, keputusan perancangan,
demonstrasi artefak melalui realisasi sistem, dan evaluasi kesesuaian hasil.
Tahapan yang digunakan mengikuti model \textit{Design Science Research
Methodology} dari Peffers dkk. \autocite{peffers2007dsrm}, sebagaimana
diringkas pada Gambar~\ref{fig:metodologi-dsr}. Pada tugas akhir ini, tahap
perancangan menghasilkan arsitektur sistem, arsitektur data, alur informasi,
dan pembagian fungsi antarkomponen, sedangkan tahap evaluasi menilai alur data
utama, pembentukan riwayat kontainer, dan integrasi antarkomponen.

\begin{figure}[ht]
  \centering
  \resizebox{0.96\textwidth}{!}{%
  \begin{tikzpicture}[
      font=\scriptsize,
      dsrstep/.style={rectangle, rounded corners, draw=black, align=center, minimum width=2.55cm, minimum height=0.84cm},
      flow/.style={-{Latex[length=1.8mm]}, thick}
    ]
    \node[dsrstep] (problem) at (0,0) {Identifikasi\\Masalah};
    \node[dsrstep] (objective) at (3.0,0) {Penetapan\\Tujuan};
    \node[dsrstep] (design) at (6.0,0) {Perancangan\\Artefak};
    \node[dsrstep] (demo) at (9.0,0) {Demonstrasi\\Sistem};
    \node[dsrstep] (evaluation) at (12.0,0) {Evaluasi\\Kesesuaian};
    \node[dsrstep] (communication) at (15.0,0) {Komunikasi\\Hasil};

    \draw[flow] (problem) -- (objective);
    \draw[flow] (objective) -- (design);
    \draw[flow] (design) -- (demo);
    \draw[flow] (demo) -- (evaluation);
    \draw[flow] (evaluation) -- (communication);
  \end{tikzpicture}}
  \caption{Alur Metodologi \textit{Design Science Research} pada Tugas Akhir}
  \label{fig:metodologi-dsr}
\end{figure}

\FloatBarrier

\section{Sistematika Penulisan}

Laporan tugas akhir ini disusun dalam tujuh bab utama dengan alur pembahasan
sebagai berikut. Bab I memuat pendahuluan dan batas kontribusi. Bab II
membahas landasan teori, regulasi, standar industri, dan penelitian terkait.
Bab III menjelaskan kondisi sistem saat ini, kebutuhan sistem, alternatif
solusi, dan dasar pemilihan pendekatan. Bab IV menyajikan rancangan arsitektur
sistem, arsitektur data, alur informasi, pembagian fungsi antarkomponen, dan
prinsip integrasi. Bab V menjelaskan realisasi sistem sebagai artefak
demonstrasi. Bab VI membahas metode, hasil, pembahasan, serta keterbatasan
evaluasi. Bab VII berisi kesimpulan dan saran pengembangan.
